\documentclass{article}
\usepackage[utf8]{inputenc}
\usepackage[portuguese]{babel}
\usepackage{graphicx}
\usepackage{booktabs}
\usepackage{amsmath}
\usepackage{hyperref}

\title{Relatório de Análise de Resultados: Q-Learning no Ambiente Taxi-v4}
\author{Gemini CLI Analysis Agent}
\date{Maio de 2026}

\begin{document}

\maketitle

\section{Introdução}
Este documento descreve os resultados obtidos a partir de uma série de experimentos e análises realizados sobre um agente de Aprendizado por Reforço treinado no ambiente \texttt{Taxi-v4} utilizando o algoritmo Q-Learning. O objetivo é fornecer dados quantitativos e qualitativos para a composição da seção de resultados de um artigo científico.

\section{Metodologia e Configurações Base}
Os experimentos foram conduzidos utilizando a biblioteca \texttt{Gymnasium}. As configurações padrão (baseline) utilizadas no treinamento foram:
\begin{itemize}
    \item \textbf{Alpha ($\alpha$):} 0.1 (Taxa de aprendizado)
    \item \textbf{Gamma ($\gamma$):} 0.6 (Fator de desconto)
    \item \textbf{Epsilon ($\epsilon$):} Inicial em 1.0, mínimo em 0.01, com decaimento de 0.001 por episódio.
    \item \textbf{Episódios de Treino:} 100.000 (para análise base) e 30.000 a 40.000 (para estudos de sensibilidade).
    \item \textbf{Semente Aleatória:} 42 (base).
\end{itemize}

\section{Resultados da Análise}

\subsection{Curva de Aprendizado e Convergência}
A análise da recompensa total por episódio revelou que o agente apresenta uma convergência clara. Inicialmente, as recompensas são altamente negativas (devido a movimentos aleatórios e penalidades), estabilizando-se em valores positivos conforme o $\epsilon$ decai.
\begin{itemize}
    \item \textbf{Observação:} A média móvel (janela de 1000) mostra que a convergência significativa ocorre por volta do episódio 15.000 a 20.000.
    \item \textbf{Arquivo Gerado:} \texttt{output/learning\_curve.png}
\end{itemize}

\subsection{Eficiência e Adesão às Regras}
\begin{itemize}
    \item \textbf{Passos por Episódio:} O número de passos médio cai drasticamente de 99 (limite máximo) para aproximadamente 13-15 passos, indicando que o agente aprendeu caminhos otimizados.
    \item \textbf{Penalidades:} As penalidades por ações ilegais (pickup/drop-off incorretos) reduzem-se a quase zero após a fase inicial de exploração, demonstrando que o agente compreendeu as restrições do ambiente.
    \item \textbf{Arquivos Gerados:} \texttt{output/steps\_per\_episode.png}, \texttt{output/penalties\_per\_episode.png}
\end{itemize}

\subsection{Sensibilidade de Hiperparâmetros ($\alpha$ e $\gamma$)}
Foram testadas variações para entender o impacto no aprendizado:
\begin{itemize}
    \item \textbf{Alpha ($\alpha$):} Testados valores [0.01, 0.1, 0.5]. O valor 0.01 mostrou-se excessivamente lento, não atingindo o pico de performance em 30k episódios. Os valores 0.1 e 0.5 convergiram de forma similar, com 0.5 sendo ligeiramente mais agressivo no início.
    \item \textbf{Gamma ($\gamma$):} Testados valores [0.1, 0.6, 0.99]. O ambiente Taxi, por ter recompensas imediatas claras, não demonstrou extrema sensibilidade ao Gamma, mas valores intermediários (0.6) apresentaram uma estabilidade ligeiramente superior.
    \item \textbf{Arquivos Gerados:} \texttt{output/sensitivity\_alpha.png}, \texttt{output/sensitivity\_gamma.png}
\end{itemize}

\subsection{Dinâmica de Exploração ($\epsilon$-decay)}
A correlação entre a queda de $\epsilon$ e o aumento da recompensa é direta. 
\begin{itemize}
    \item \textbf{Decaimento Rápido (0.005):} Atinge a performance máxima muito cedo, o que é eficiente para este ambiente estável.
    \item \textbf{Decaimento Lento (0.0001):} Mantém o agente explorando por mais tempo que o necessário, atrasando a estabilização da recompensa média.
    \item \textbf{Arquivo Gerado:} \texttt{output/sensitivity\_epsilon\_decay.png}
\end{itemize}

\subsection{Robustez e Generalização}
Para garantir que os resultados não foram fruto de sorte estatística, o modelo foi treinado com 5 sementes diferentes (42, 10, 2024, 999, 12345).
\begin{itemize}
    \item \textbf{Resultado:} O desvio padrão entre as corridas foi mínimo, comprovando a robustez do algoritmo Q-Learning para este problema. Todas as sementes convergiram para patamares de recompensa similares.
    \item \textbf{Arquivo Gerado:} \texttt{output/model\_robustness.png}
\end{itemize}

\section{Discussão}
O Q-Learning mostrou-se uma solução extremamente eficaz para o problema do Taxi-v4. A escolha dos hiperparâmetros padrão foi validada pelas análises de sensibilidade. Observou-se que a fase de exploração ($\epsilon$) é o fator mais crítico para a velocidade de aprendizado inicial, enquanto o Alpha define a estabilidade da convergência final. A eficiência (passos) e a segurança (penalidades) convergiram simultaneamente com a recompensa, indicando um aprendizado coerente da política ótima.

\end{document}
